% =========================================================================
% Chapter: Background
% Dissertation: Computer Vision for Urban Waste and Infrastructure
%               Monitoring in Malta (University of Malta, 2026)
% Written against repository commit 1b2c918a (inspected 11 July 2026);
% revised against commit 2d70be27 (inspected 16 July 2026) to add the
% LingBot-Vision representation family; re-verified against commit 0c344db
% (18 July 2026) — no factual changes required (chapter states no counts).
% Citation style: natbib author-year (\citet / \citep).
% =========================================================================

\chapter{Background}
\label{ch:background}

This chapter establishes the conceptual and terminological foundations on which the dissertation rests. It introduces municipal monitoring as an application domain, situates it within the Maltese institutional context, defines the computer-vision tasks that the dissertation draws upon, and describes the two datasets around which the experimental work is organised: the Maltese Domestic Waste Dataset (MDWD) and the Maltese Traffic Sign Dataset (MTSD). It then presents, at a conceptual level, the three families of technical approaches compared in this work---supervised object detection, prompt-based foundation models, and representation-learning attribute classification---together with the deployment and evaluation concepts needed to interpret the later chapters. Detailed comparisons of individual published systems are deferred to Chapter~\ref{ch:literature-review}, and the dissertation's own experimental procedures are described in the Methodology chapter.

\section{Urban and Municipal Monitoring}
\label{sec:bg-municipal}

Municipalities are responsible for a continuous cycle of inspection, intervention and maintenance across public space: streets must be kept clean, domestic waste must be collected reliably, and road furniture such as traffic signs must remain visible, legible and structurally sound. In practice much of this work depends on knowing, at any given time, \emph{where} a problem exists---a pile of uncollected refuse, an illegally dumped item, a corroded or vandalised sign---and this situational awareness has traditionally been obtained through manual inspection, scheduled patrols and citizen complaints. Manual inspection is labour-intensive, covers limited geographical area per staff-hour, and is necessarily intermittent: a street inspected weekly can accumulate and retain problems for days before they are recorded. Missed or delayed observations have concrete consequences. Uncollected or incorrectly presented waste attracts pests, obstructs narrow pavements, and degrades the perceived cleanliness of a locality; a faded, bent or obscured traffic sign degrades the information available to drivers and pedestrians precisely at locations where that information is safety-relevant. Timely, georeferenced visual evidence therefore has direct operational value: it allows maintenance and enforcement resources to be prioritised by observed condition rather than by fixed schedules, and it creates an auditable record of the state of public assets. Automated visual monitoring should be understood in this support role---as a way of extending the coverage and frequency of observation available to municipal staff---rather than as a wholesale replacement for human judgement about intervention and enforcement \citep{chen2019deep,merolla2025improving}.

\subsection{The Maltese Context}
\label{sec:bg-malta}

Malta presents an unusually concentrated version of the urban-monitoring problem. It is among the most densely populated states in Europe, and its towns are characterised by narrow streets, continuous terraced frontages, on-street parking, and high visual clutter in which waste bags, bins, signage, vehicles and street furniture compete for limited space. Kerbside waste is presented directly on pavements rather than in dedicated container areas, so it is routinely partially occluded by parked cars and pedestrians. Coastal exposure accelerates the corrosion and fading of metal signage, and the signage stock itself is heterogeneous, mixing modern reflective signs with older designs, bilingual street nameplates, tourist-oriented signage and locally specific furniture such as convex blind-spot mirrors. Illumination varies sharply between bright Mediterranean sunlight, deep shadow in narrow streets, and low-light conditions. These characteristics---clutter, occlusion, small object size, illumination variability and ageing infrastructure---are exactly the conditions under which computer-vision systems trained on foreign benchmark data are least reliable, which motivates the collection of Malta-specific datasets in this dissertation.

Institutional responsibility for the two monitored domains is distributed across several bodies, and the division of labour matters when interpreting the datasets. Household waste separation is mandatory in Malta and follows a national colour convention: organic waste is presented in white bags, co-mingled recyclables (paper, plastic and metal) in grey or green bags, and residual mixed waste in black bags, with glass collected separately on designated days; kerbside collection is organised through local and regional councils under a nationally harmonised schedule \citep{wastecollection2026web,era2023web}. WasteServ Malta operates the downstream treatment and disposal facilities to which collected streams are delivered \citep{wasteserv2026web}. Public cleansing---the cleaning of arterial roads, removal of illegally dumped waste, beach cleaning and related upkeep---is the remit of the Cleansing and Maintenance Division (CMD), which provides these services to central government and to local councils \citep{cmd2026web}. Road infrastructure responsibilities are likewise shared: Transport Malta acts as the transport regulator and issues the technical guidelines governing traffic signs \citep{transportmalta2021guidelines}; Infrastructure Malta has progressively assumed responsibility for the construction and maintenance of arterial and distributor roads, publishing the engineering specifications under which signs are procured and installed \citep{infrastructuremalta2020spec}; and local councils retain maintenance responsibility for residential roads under the Local Councils Act \citep{localcouncilsact1993}. Responsibility for any particular sign or waste presentation therefore depends on road classification, location and service type, and a monitoring system that produces georeferenced evidence is useful to several of these actors simultaneously rather than to a single authority.

\section{Computer-Vision Tasks Relevant to this Dissertation}
\label{sec:bg-tasks}

The dissertation draws on several distinct computer-vision task formulations, and the differences between them determine what each family of models can and cannot be asked to do. \emph{Image classification} assigns one or more labels to an entire image or to a cropped region, without indicating where in the image the evidence lies \citep{krizhevsky2012imagenet}. \emph{Object detection} predicts a set of bounding boxes, each with a category label and, in most supervised detectors, a confidence score, thereby answering both \emph{what} and \emph{where} \citep{everingham2010pascal,lin2014microsoft}. \emph{Semantic segmentation} assigns a class to every pixel; \emph{instance segmentation} additionally separates individual object instances, producing per-object masks rather than boxes \citep{he2017mask}. \emph{Object localisation} is used here as a general term for producing spatial evidence (boxes or points) for indicated objects, whether or not calibrated scores accompany it. \emph{Open-vocabulary detection} relaxes the assumption of a fixed training label set: categories are specified at inference time as natural-language text, so a model can in principle be asked for classes it never saw as explicit training labels \citep{zareian2021open}. \emph{Prompt-based segmentation} produces masks in response to user prompts such as points, boxes or text \citep{kirillov2023segment}. \emph{Visual question answering and multimodal reasoning} treat the image as context for a language model that generates free-form text, which may include structured outputs such as coordinates. Finally, \emph{attribute classification} predicts properties of an object that has already been localised---in this dissertation, properties of a traffic sign crop such as its physical condition.

Several classification distinctions recur later and are worth fixing precisely. A \emph{multi-class} problem selects exactly one label from a set of mutually exclusive classes; a \emph{multi-label} problem allows several labels to hold simultaneously. \emph{Multi-attribute} classification, as used here, denotes several parallel classification problems over the same object, each attribute having its own mutually exclusive class vocabulary: every MTSD sign receives exactly one viewing-angle label, one mounting label, one condition label and one shape label. Such a problem can be solved either by separate single-head models, one per attribute, or by a \emph{multi-head} model in which a shared backbone feeds several task-specific output heads; the trade-offs between these designs are discussed in Section~\ref{sec:bg-multihead}.

\section{The MDWD Domain: Kerbside Domestic Waste}
\label{sec:bg-mdwd}

The Maltese Domestic Waste Dataset (MDWD) is a street-level object-detection dataset representing domestic waste as it is actually presented for kerbside collection in Malta. Its five classes are operational categories: they encode how waste appears and is presented for collection, not a laboratory taxonomy of materials. This distinction matters because a sealed, opaque bag does not reveal its contents; annotation is necessarily based on visible presentation---bag colour, form, context and placement---and the class definitions below should be read as visual-operational conventions rather than as claims about the true material composition of any individual bag.

\paragraph{Mixed Waste.} This class represents residual general domestic waste, most commonly opaque black refuse bags and tied domestic sacks that are not visibly assigned to a separated recycling or organic stream, consistent with the national black-bag convention for residual waste \citep{wastecollection2026web}. Visually the class spans single bags, groups of bags placed together at a collection point, and deformed or partially collapsed bags. Dark, wet or reflective plastic is easily confused with shadows, tyres and other dark street objects, and clusters of touching bags challenge the separation of individual instances.

\paragraph{Orange CMD.} This class covers orange waste bags and receptacles of the type associated with municipal cleansing activity carried out by or for the Cleansing and Maintenance Division \citep{cmd2026web}. The identifying evidence is the distinctive orange bag colour in a street-cleansing context. Because orange objects of many other kinds appear in streets, colour alone is not a sufficient criterion; placement, bag form and context distinguish the class, and the annotation policy treats it as an operational category of the dataset rather than a claim about bag contents.

\paragraph{Organic Waste.} This class represents separately presented biodegradable household waste, which in Malta is presented in white bags under the national separation scheme \citep{wastecollection2026web,era2023web}. The class denotes the separated organic collection stream, not arbitrary food-like material visible in the street. Visual ambiguity arises from overfilled or stretched bags, semi-transparent bags whose contents show through, soiled bags, poor illumination, and non-standard bags used by some households.

\paragraph{Recyclable Material.} This class represents the separated recyclable stream, presented in grey or green bags of co-mingled paper, plastic and metal \citep{wastecollection2026web}, together with visibly recyclable presentations such as flattened cardboard placed for collection. The annotation policy follows collection presentation rather than material inference: an object is not labelled recyclable merely because its material could in principle be recycled, and mixed piles in which recyclable items lie among residual waste are a recognised source of inter-class confusion.

\paragraph{Other Waste.} This is the residual operational category for waste that is clearly relevant to the monitoring task but does not satisfy the definitions above---for example bulky items, loose discarded objects and containers, and other non-standard kerbside presentations. Retaining an explicit residual class keeps the remaining four classes semantically clean while still allowing the detector to flag waste presence that would otherwise be unannotated.

Taken together, the taxonomy is deliberately aligned with the collection streams that Maltese municipal actors operate, so that per-class detections translate directly into operationally meaningful observations (for instance, organic bags presented on a non-organic collection day). The principal sources of inter-class confusion are colour degradation under poor lighting, occlusion by vehicles, physical contact between bags of different streams, and the inherently residual definitions of Mixed Waste and Other Waste.

\section{The MTSD Domain: Traffic-Sign and Roadside Infrastructure}
\label{sec:bg-mtsd}

The Maltese Traffic Sign Dataset (MTSD) is a street-level dataset of Maltese traffic signs and closely related roadside infrastructure, annotated with bounding boxes over twelve object classes and, for every box, four per-sign attributes (Section~\ref{sec:bg-attributes}). The twelve classes were selected as recurrent, driver-relevant and operationally significant categories within the study's collection scope, covering regulatory, warning, navigational and supporting infrastructure; they are not claimed to be a statistically complete inventory of Maltese signage. Annotators were instructed to box every visible object belonging to the taxonomy and to exclude private signage, commercial notices and advertisements.

Six classes carry conventional traffic-control semantics. \emph{Pedestrian Crossing} covers signs warning of or marking pedestrian crossings, which are safety-critical because they mediate pedestrian priority and driver speed adaptation. \emph{Stop Sign} covers the octagonal regulatory stop sign, among the most safety-significant and visually distinctive signs in any inventory. \emph{No Entry (One Way)} covers the regulatory no-entry sign as deployed in Malta's dense one-way street networks; the dataset label reflects the local pairing of no-entry signage with one-way circulation rather than an assertion that the two concepts are legally identical. \emph{Roundabout Ahead} covers the warning sign for an approaching roundabout, relevant to speed and yielding behaviour on Malta's roundabout-dense arterial network. \emph{No Through Road (T-Junction)} is a combined dataset category for the locally recurring signage indicating dead-end and T-configuration road layouts; the label names the dataset category, and its precise membership is defined by the annotation instructions rather than by a single legal sign definition. \emph{Blind-Spot Mirror (Convex Mirror)} is not a traffic sign in the regulatory sense but roadside visibility infrastructure: convex mirrors installed at restricted-visibility junctions, bends and entrances, which are common in Malta's narrow streets and are maintained as road furniture.

Three classes carry navigational or informational semantics. \emph{Street Sign} covers street nameplates, which in Malta are frequently wall-mounted, bilingual and of heterogeneous age and design; they matter for navigation, addressing, emergency response and asset management. \emph{Directional Sign} covers signs directing road users towards destinations, routes and facilities. \emph{Tourist Sign} covers signage indicating tourist attractions, heritage locations and visitor facilities, a locally prominent category given Malta's tourism economy. \emph{Auxiliary Sign} covers supplementary plates---typically white with a dark border, carrying text, symbols, times, distances, vehicle restrictions or exceptions---whose defining property is functional rather than chromatic: an auxiliary sign qualifies or extends the meaning of another sign rather than carrying free-standing meaning of its own.

The final two classes handle irreducible uncertainty explicitly. \emph{Back-Unknown} is applied when the rear of a sign is visible, the object is confidently recognisable as traffic-sign infrastructure, but the front-facing category cannot be determined. \emph{Other-Unknown} is applied when an object clearly belongs to public traffic or roadside sign infrastructure but matches none of the predefined semantic classes. The two are distinguished by the source of uncertainty: Back-Unknown encodes rear-facing visibility with unknown semantics, whereas Other-Unknown encodes an unmatched infrastructure type. Retaining these classes is preferable to forcing unsupported semantic guesses, both for annotation honesty and because rear-facing signs remain legitimate detection targets for inventory purposes.

\subsection{The MTSD Attribute Taxonomy}
\label{sec:bg-attributes}

Every annotated MTSD box carries four attributes, each with a small, mutually exclusive class vocabulary. The attributes describe the sign as observed, and together they turn a detection inventory into a condition-assessment resource.

\emph{Viewing angle} (Front, Side, Back) records the orientation of the sign relative to the capturing camera---not relative to geographic north, the carriageway or traffic flow. Front means the sign face is substantially visible and its content normally interpretable; Side means a strongly oblique or edge-on view with much-reduced visible face area; Back means the rear surface dominates the view.

\emph{Mounting type} (Pole-Mounted, Wall-Mounted) records the supporting structure: Pole-Mounted covers signs supported by poles, posts or similar ground-rising structures, including shared and utility poles, while Wall-Mounted covers signs fixed to walls, fa\c{c}ades or wall-attached brackets, including signs projecting outward from a wall. Wall mounting is disproportionately common in Malta because of its narrow pavements and continuous building frontages.

\emph{Physical condition} (Good, Weathered, Heavily Damaged) is an ordinal assessment of deterioration. A \emph{Good} sign is structurally intact, clearly recognisable and legible, essentially free of rust, substantial fading, bending, cracks or missing sections, and at most minimally affected by dirt, stickers or minor marks; very small cosmetic imperfections remain compatible with Good where they do not meaningfully affect presentation or legibility. A \emph{Weathered} sign shows clearly visible ageing---localised rust, moderate fading or discolouration, paint wear, sticker residue, minor dents, superficial corrosion or moderate dirt---while remaining functionally recognisable and substantially legible; deterioration is evident but does not amount to structural or functional failure. A \emph{Heavily Damaged} sign exhibits critical deterioration beyond cosmetic weathering: extensive corrosion, large missing or broken sections, holes, severe bending or deformation, major cracks, substantial peeling, heavy fading or obscuration, severe vandalism, or loss of structural integrity, such that readability, recognisability or safe function is materially affected. The boundary between Weathered and Heavily Damaged is thus drawn at the transition from visible-but-tolerable ageing to damage that would justify prioritised replacement.

\emph{Sign shape} (Circular, Quadrangle, Triangular, Octagonal, Pentagon) records the geometric outline of the sign face, with Quadrangle serving as the dataset's umbrella category for all four-sided signs, whether square or rectangular. Shape is strongly correlated with regulatory function in European signage conventions (circular prohibitions and obligations, triangular warnings, octagonal stop), so it provides an annotation-level consistency check as well as a classification target.

\section{Object-Detector Families}
\label{sec:bg-detectors}

Modern object detectors form a small number of architectural families. \emph{Two-stage} detectors first generate class-agnostic region proposals and then classify and refine each proposal; Faster R-CNN is the canonical example \citep{ren2015faster}. \emph{One-stage} detectors instead predict categories and box offsets directly over a dense grid of image locations in a single pass, trading some accuracy---historically---for simplicity and speed \citep{redmon2016you,liu2016ssd,lin2017focal}. Within both families, \emph{anchor-based} prediction regresses offsets from predefined reference boxes, whereas \emph{anchor-free} prediction regresses box geometry directly from feature-map locations. Convolutional backbones extract hierarchical features \citep{he2016deep}, usually fused across scales by a feature pyramid so that small objects are detected from high-resolution levels \citep{lin2017feature}. Attention mechanisms and Transformer encoder--decoder blocks \citep{dosovitskiy2021image} entered detection through DETR-style architectures, which recast detection as \emph{set prediction}: a fixed set of learned queries attends to image features, and training uses bipartite \emph{matching} so that each ground-truth object is assigned to exactly one prediction \citep{carion2020end}. This one-to-one formulation removes the need for \emph{non-maximum suppression} (NMS), the post-processing step by which dense detectors discard near-duplicate overlapping predictions \citep{neubeck2006efficient}. Detectors are therefore usefully classified along two further axes: NMS-based versus end-to-end (NMS-free) prediction, and training from random initialisation versus transfer learning from pretrained weights. The supervised track of this dissertation spans these axes with three recent YOLO generations (YOLO11, YOLO12 and YOLO26) and one real-time detection transformer (RF-DETR); their published characteristics are reviewed critically in Chapter~\ref{ch:literature-review}.

\section{Foundation Models and Prompt-Based Perception}
\label{sec:bg-foundation}

A \emph{foundation model} is a large model pretrained on broad data that is adapted---often without weight updates---to many downstream tasks \citep{bommasani2021opportunities}. Visual foundation models specialise this idea to images and video. Several distinct capabilities are relevant here and must not be conflated. \emph{Promptable segmentation} models produce masks for objects indicated by point, box or text prompts \citep{kirillov2023segment}. \emph{Open-vocabulary detectors} align region features with text embeddings so that arbitrary category names can be queried at inference time \citep{zareian2021open,li2022grounded}. \emph{Vision-language models} (VLMs) couple a visual encoder to a language model and generate text conditioned on images \citep{radford2021learning,bai2025qwen25vl}; when that text is constrained to contain coordinates, the VLM performs \emph{grounded localisation}. \emph{Multimodal large language models} extend this with large-scale language reasoning. \emph{World foundation models} are models of physical environments and their dynamics, developed primarily for physical AI applications such as robotics and autonomous systems \citep{ha2018world,nvidia2025cosmosplatform}; reasoning-oriented members of such ecosystems, such as NVIDIA's Cosmos Reason models, are VLMs post-trained for physically grounded reasoning \citep{nvidia2025cosmosreason1}. \emph{Zero-shot} inference applies a pretrained model to a new task without task-specific training; \emph{few-shot} inference additionally supplies a handful of examples, typically in the prompt \citep{brown2020language}.

Prompt type is an equally important distinction: point, box and mask prompts indicate \emph{where} without naming a category, whereas text prompts name a category without indicating location. Natural-language prompting does not guarantee correct class grounding: the mapping from a phrase such as ``organic-waste bag'' to pixels depends on the model's training distribution, and synonym choice, phrasing, ambiguity and scene context can all change the output. Finally, systems differ in output form---mask-producing versus box-producing---and in whether they emit calibrated per-detection confidence scores; generative VLMs typically produce boxes as text tokens with no associated score, a difference with direct consequences for evaluation (Section~\ref{sec:bg-evaluation}).

\section{Representation Learning and Transfer}
\label{sec:bg-representation}

When labelled domain data are scarce---as for any newly collected municipal dataset---most of a model's visual competence must come from pretraining. \emph{Supervised pretraining} learns features from large labelled corpora such as ImageNet \citep{deng2009imagenet,krizhevsky2012imagenet}. \emph{Self-supervised learning} (SSL) removes the label requirement by constructing pretext objectives from the data itself: contrastive methods pull augmented views of the same image together \citep{chen2020simple,he2020momentum}, self-distillation methods train a student to match a momentum teacher \citep{grill2020bootstrap,caron2021emerging}, and masked-image modelling reconstructs hidden patches \citep{he2022masked}. \emph{Joint-embedding predictive architectures} (JEPAs) predict the latent representation of masked content rather than its pixels, an approach argued to favour semantic abstraction over low-level detail \citep{lecun2022path,assran2023self}, and extend naturally to video, where prediction across time encourages spatiotemporal structure \citep{bardes2024revisiting}. Pretrained representations are transferred to a downstream task along a spectrum of adaptation strategies: \emph{linear probing} freezes the backbone and trains only a shallow head, measuring representation quality at minimal cost; \emph{parameter-efficient fine-tuning}, exemplified by low-rank adaptation (LoRA), inserts small trainable low-rank matrices into selected layers while freezing the original weights \citep{hu2022lora,houlsby2019parameter}; and \emph{full fine-tuning} updates all parameters, maximising adaptation capacity at the highest computational and overfitting risk \citep{yosinski2014transferable,kornblith2019better}. The attribute-classification track of this dissertation compares these strategies across four pretrained representation families that differ in their pretraining objective: a general-purpose self-supervised image transformer (DINOv3), a video-pretrained JEPA (V-JEPA~2.1), a supervised modern convolutional network (ConvNeXt), and a boundary-centric self-supervised transformer trained for dense spatial perception (LingBot-Vision). Their published evidence is reviewed in Chapter~\ref{ch:literature-review}.

\section{Multi-Head and Single-Head Classification}
\label{sec:bg-multihead}

Predicting the four MTSD attributes poses an architectural choice. A \emph{shared-backbone multi-head} classifier computes one representation of the sign crop and feeds it to four independent output heads, one per attribute. This design is a form of hard parameter sharing in multi-task learning \citep{caruana1997multitask,ruder2017overview}. Its potential benefits are substantial: the four attributes describe the same physical object, so features useful for one (for example, edge geometry for shape) may inform another (for example, viewing angle); backbone computation is performed once rather than four times, reducing training cost, inference latency and deployment footprint; and a single model is simpler to version, monitor and update in production. Its risks are equally documented: \emph{negative transfer}, in which jointly trained tasks degrade one another; conflicting gradients pulling shared parameters in incompatible directions \citep{yu2020gradient}; unequal task difficulty and class imbalance allowing one attribute to dominate the shared representation; different convergence rates across heads; and sensitivity to how the per-head losses are weighted \citep{kendall2018multi,chen2018gradnorm}. The alternative---four separate single-head models---eliminates inter-task interference and lets each model be tuned independently, at the price of quadrupled computation, storage and maintenance, and without any opportunity for beneficial feature sharing. Which regime prevails is an empirical question that depends on task relatedness \citep{standley2020tasks}; the dissertation adopts the shared-backbone multi-head design for all attribute experiments, and the choice is justified operationally in the Methodology chapter.

\section{Edge Computing and Real-Time Monitoring}
\label{sec:bg-edge}

Municipal visual monitoring ultimately runs somewhere, and the choice of \emph{where} shapes which models are viable. At one extreme, imagery is transmitted to centralised cloud or on-premises servers with large GPUs; at the other, inference runs \emph{on-device} on the capturing hardware (a smartphone, dashcam or embedded computer in an inspection vehicle); between them lie near-edge gateways and hybrid designs in which lightweight edge models filter or pre-screen data for heavier server-side analysis \citep{shi2016edge,satyanarayanan2017emergence,chen2019deep}. The trade-offs are well documented \citep{murshed2021machine}. Server-side inference offers effectively unbounded compute and centralised model updates, but requires transmitting street imagery over mobile networks---raising bandwidth cost, latency, and exposure of personal data such as faces and number plates to transfer and storage, with corresponding obligations under the General Data Protection Regulation \citep{gdpr2016}. Edge inference keeps imagery local, tolerates intermittent connectivity (relevant in dense urban canyons), and bounds operating cost, but is constrained by device memory, thermal envelopes, energy budgets and the difficulty of updating models in the field.

``Real time'' is an application-dependent property, not a model property. Model inference latency is only one term in end-to-end pipeline latency, which also includes camera capture, preprocessing, any network transfer, post-processing such as duplicate suppression, and result handling; throughput (frames per second) additionally depends on batch size and hardware, and a system quoted at high FPS on a server GPU with large batches may behave very differently on a CPU with batch size one \citep{chen2019deep}. For municipal monitoring from moving vehicles, dashcams, smartphones or fixed roadside cameras, the binding requirement is usually that processing keeps pace with the capture rate at acceptable energy cost, not that any single frame is processed within milliseconds. These considerations explain why speed--accuracy--memory trade-offs, rather than accuracy alone, structure the model comparisons in this dissertation: a smaller model that runs on affordable hardware at the required rate may be operationally preferable to a marginally more accurate model that does not, and techniques such as quantisation and pruning exist precisely to move models along this trade-off curve \citep{han2016deep,jacob2018quantization,howard2017mobilenets}.

\section{Evaluation Terminology}
\label{sec:bg-evaluation}

The later chapters assume standard evaluation vocabulary, summarised here generically; the dissertation's specific protocol is defined in the Methodology chapter. For object detection, a predicted box is matched to a ground-truth box by \emph{intersection over union} (IoU), the ratio of overlap area to union area. Given an IoU threshold and a \emph{confidence threshold}, each prediction is a true positive (matched), false positive (unmatched prediction) or contributes a false negative (unmatched ground truth). \emph{Precision} is the fraction of predictions that are correct, \emph{recall} the fraction of ground-truth objects found, and the \emph{F1 score} their harmonic mean. Sweeping the confidence threshold traces a precision--recall curve whose area is the \emph{average precision} (AP) for a class; the mean over classes at IoU~0.50 is mAP@50, and averaging additionally over IoU thresholds from 0.50 to 0.95 yields the stricter mAP@50:95 popularised by the COCO benchmark \citep{lin2014microsoft,everingham2010pascal,padilla2021comparative}. Per-class results matter whenever class frequencies are imbalanced, as they are in both datasets studied here. Importantly, AP presupposes rankable confidence scores; detectors that emit unscored boxes can only be evaluated at a single operating point (precision, recall, F1), which constrains cross-paradigm comparison \citep{guo2017calibration}.

For classification, \emph{accuracy} is the fraction of correct predictions; per-class precision, recall and F1 are aggregated either as \emph{macro} averages, which weight all classes equally and therefore expose minority-class performance, or as \emph{weighted} averages, which follow class frequency \citep{sokolova2009systematic}. A \emph{confusion matrix} tabulates predicted against true classes and reveals systematic error structure. Under class imbalance---severe for the MTSD condition attribute, where damaged signs are rare---accuracy is a misleading summary and macro-F1 is the more informative headline metric. For deployment comparison, the relevant quantities are parameter count, checkpoint size on disk, inference latency and throughput on stated hardware, peak GPU memory, CPU-only feasibility, and, where measurable, energy use \citep{murshed2021machine}.
